\section{Vehicles}

\subsection{Biovore}

The Biovore is a squat and bloated creature, but this makes it no less dangerous. Its massive body constantly produces Spore-Mines: living bombs that explode in a spray of acid and razor-sharp bone fragments.

Biovores advance ponderously across the battlefield, using the bony protrusions on their forelimbs to anchor themselves to the ground before firing. With a powerful muscular spasm, the Biovore launches its Spore-Mines into the air. Each mine contains a gas-filled bladder that slows its descent.

Even if the shot misses its intended target, the Spore-Mine remains dangerous. Its rudimentary intelligence prevents it from detonating upon impact, instead causing it to drift towards nearby prey before exploding. Each Biovore fires one Spore-Mine per turn.

The Biovore has the \textbf{Walker} and \textbf{Slave (Nest)} abilities. Its Spore-Mines have the \textbf{Bio-Toxin}, \textbf{Artillery}, \textbf{Battery}, \textbf{Template (12 cm)}, \textbf{Spore-Mine}, and \textbf{Reduces Cover ($-1$)} abilities.

\begin{tblr}{
    width=\textwidth,
    colspec={
        Q[c,m,1.2]
        Q[c,m,1.2]
        Q[c,m,0.7]
        Q[c,m,0.8]
        X[c,m,2.7]
        Q[c,m,1.1]
        Q[c,m,0.9]
        Q[c,m,1.2]
        Q[c,m,0.7]
    },
    row{1}={bg=black, fg=white, font=\bfseries},
    cells={valign=m},
    hlines,
    vlines
}
Movement & Save & AF & Morale & Weapon & Range & Dice & To-Hit & AP \\
15 cm & 4+ & +0 & 6 & Spore-Mines & 90 cm & Template & 4+ & $-1$
\end{tblr}

\unitdivider


\subsection{Dactylis}

Dactylis are creatures specialised in long-range fire, delivering deadly barrages upon the prey of the Hive. Their Bile Pods can eat through even the strongest armour and plasteel, seeping through walls and defensive barriers.

The Dactylis has the \textbf{Walker} and \textbf{Slave (Nest)} abilities. Its Bile Pods have the \textbf{Artillery} ability. The Bile Pods -- Spores firing mode also has the \textbf{Reduces Cover ($-2$)} ability.

\begin{tblr}{
    width=\textwidth,
    colspec={
        Q[c,m,1.2]
        Q[c,m,1.2]
        Q[c,m,0.7]
        Q[c,m,0.8]
        X[c,m,2.7]
        Q[c,m,1.1]
        Q[c,m,0.9]
        Q[c,m,1.2]
        Q[c,m,0.7]
    },
    row{1}={bg=black, fg=white, font=\bfseries},
    cells={valign=m},
    hlines,
    vlines
}
Movement & Save & AF & Morale & Weapon & Range & Dice & To-Hit & AP \\

\SetCell[r=2]{c,m} 15 cm &
\SetCell[r=2]{c,m} 3+ &
\SetCell[r=2]{c,m} +1 &
\SetCell[r=2]{c,m} 6 &
Bile Pods -- Spores &
90 cm &
2 &
3+ &
0
\\

& & & & 
Bile Pods -- Bio-Acid &
90 cm &
2 &
4+ &
$-2$
\end{tblr}

\unitdivider


\subsection{Exocrine}

The Exocrine is effectively the means of transportation for the symbiotic weapon attached to its back. Although it possesses colossal physical strength, the Exocrine has a tiny brain whose capabilities are far inferior to those of the bio-weapon it carries.

It is not uncommon for the weapon to direct the Exocrine's behaviour, guiding it towards an advantageous firing position. Only when its host has become stationary can the weapon concentrate fully upon destroying its prey.

Despite the Exocrine's relative simplicity, it is perfectly adapted to its role. Batteries of these ponderous beasts can devastate armoured columns or tightly packed formations of infantry.

The Exocrine has the \textbf{Walker} and \textbf{Slave (Nest)} abilities.

\begin{tblr}{
    width=\textwidth,
    colspec={
        Q[c,m,1.2]
        Q[c,m,1.2]
        Q[c,m,0.7]
        Q[c,m,0.8]
        X[c,m,2.7]
        Q[c,m,1.1]
        Q[c,m,0.9]
        Q[c,m,1.2]
        Q[c,m,0.7]
    },
    row{1}={bg=black, fg=white, font=\bfseries},
    cells={valign=m},
    hlines,
    vlines
}
Movement & Save & AF & Morale & Weapon & Range & Dice & To-Hit & AP \\
15 cm & 3+ & +1 & 6 & Bio-Plasma Cannon & 75 cm & 2 & 4+ & $-3$
\end{tblr}

\unitdivider

\subsection{Harpy}

The Harpy is a monstrous organism capable of weaving through the skies with a speed and precision that even the most advanced conventional fighter aircraft could only dream of matching.

As it glides over the prey world, it releases a rain of living bombs and spits death from the bio-weapons fused to its forelimbs. Harpies appear during the opening stages of a Tyranid attack, working alongside Gargoyles to drive their prey from cover. Although the two species share the same purpose, they are physically very different.

The Harpy has the \textbf{Semi-Synaptic (Hunt)} and \textbf{Floater} abilities. Its Spore-Mines have the \textbf{Bombardment (1)}, \textbf{Template (7.5 cm)}, \textbf{Bio-Toxin}, \textbf{Spore-Mine}, and \textbf{Battery} abilities.

\begin{tblr}{
    width=\textwidth,
    colspec={
        Q[c,m,1.2]
        Q[c,m,1.2]
        Q[c,m,0.7]
        Q[c,m,0.8]
        X[c,m,2.7]
        Q[c,m,1.1]
        Q[c,m,0.9]
        Q[c,m,1.2]
        Q[c,m,0.7]
    },
    row{1}={bg=black, fg=white, font=\bfseries},
    cells={valign=m},
    hlines,
    vlines
}
Movement & Save & AF & Morale & Weapon & Range & Dice & To-Hit & AP \\

\SetCell[r=2]{c,m} 25 cm &
\SetCell[r=2]{c,m} 3+ &
\SetCell[r=2]{c,m} +4 &
\SetCell[r=2]{c,m} 5 &
Venom Cannon &
30 cm &
2 &
4+ &
$-1$
\\

& & & &
Spore-Mines &
Bomb &
Template &
4+ &
0
\end{tblr}

\unitdivider


\subsection{Haruspex}

The Haruspex is a ferocious beast created for the sole purpose of consuming astronomical quantities of biomass. Its insatiable appetite drives it to seek out ever more prey to devour.

Few enemies are foolish enough to stand against a Haruspex, for it can consume an entire platoon within a matter of seconds, swallowing soldiers one after another into its gaping maw without ever appearing satisfied.

Prey too large to swallow whole is simply torn into more manageable pieces by its enormous claws. Buildings are smashed open and tanks ripped apart, their occupants devoured before they can even understand what is happening.

Only those who choose to flee have any chance of escape, and then only if they can avoid the Haruspex's long, extendable tongue, which strikes like lightning to seize its victims before drawing them into its cavernous mouth.

The Haruspex has the \textbf{Slave (Devastation)} and \textbf{Walker} abilities.

\begin{tblr}{
    width=\textwidth,
    colspec={
        Q[c,m,1.2]
        Q[c,m,1.2]
        Q[c,m,0.7]
        Q[c,m,0.8]
        X[c,m,2.7]
        Q[c,m,1.1]
        Q[c,m,0.9]
        Q[c,m,1.2]
        Q[c,m,0.7]
    },
    row{1}={bg=black, fg=white, font=\bfseries},
    cells={valign=m},
    hlines,
    vlines
}
Movement & Save & AF & Morale & Weapon & Range & Dice & To-Hit & AP \\
20 cm & 2+ & +7 & 6 & Hooks & 20 cm & 2 & 5+ & 0
\end{tblr}

\unitdivider


\subsection{Malefactor}

Malefactors are enormous organic transports protected by reinforced chitin plates that provide armour comparable to that of the heaviest tanks.

Armed with Fragmentation Spines, Malefactors are formidable transports capable of providing fire support to their troops after they have disembarked.

The Malefactor has the \textbf{Slave (Devastation)}, \textbf{Attached Transport}, \textbf{Walker}, and \textbf{Transport (2)} abilities.

\begin{tblr}{
    width=\textwidth,
    colspec={
        Q[c,m,1.2]
        Q[c,m,1.2]
        Q[c,m,0.7]
        Q[c,m,0.8]
        X[c,m,2.7]
        Q[c,m,1.1]
        Q[c,m,0.9]
        Q[c,m,1.2]
        Q[c,m,0.7]
    },
    row{1}={bg=black, fg=white, font=\bfseries},
    cells={valign=m},
    hlines,
    vlines
}
Movement & Save & AF & Morale & Weapon & Range & Dice & To-Hit & AP \\
20 cm & 2+ & +5 & Attached & Fragmentation Spines & 20 cm & 2 & 5+ & 0
\end{tblr}

\unitdivider

\subsection{Neurotyrant}

Neurotyrants are masses of tentacles and chitin that act as powerful synaptic nodes for the Hive Mind. This same biological structure also allows them to float slightly above the ground.

Their unique cerebral structure acts as a focal point for the Shadow in the Warp, increasing the psychic pressure upon their enemies. Neurotyrants function as mobile command centres, providing surges of energy to nearby Tyranid creatures.

The Neurotyrant has the \textbf{Leader}, \textbf{Synapse (30 cm)}, \textbf{Dark Presence ($-1$/12 cm)}, and \textbf{Walker} abilities. Its Energy Torrent has the \textbf{Reduces Cover ($-3$)} ability.

\begin{tblr}{
    width=\textwidth,
    colspec={
        Q[c,m,1.2]
        Q[c,m,1.2]
        Q[c,m,0.7]
        Q[c,m,0.8]
        X[c,m,2.7]
        Q[c,m,1.1]
        Q[c,m,0.9]
        Q[c,m,1.2]
        Q[c,m,0.7]
    },
    row{1}={bg=black, fg=white, font=\bfseries},
    cells={valign=m},
    hlines,
    vlines
}
Movement & Save & AF & Morale & Weapon & Range & Dice & To-Hit & AP \\
20 cm & 3+ & +2 & -- & Energy Torrent & 20 cm & 2 & 4+ & 0
\end{tblr}

\unitdivider


\subsection{Pyrovore}

The role of a Pyrovore is to predigest biomass. Its jaws drip with an acid capable of reducing flesh, metal, and even stone into a pulp that can easily be absorbed by other Tyranid organisms.

However, the Pyrovore's most effective weapon is the dorsal cannon capable of spitting great spheres of flame.

The Pyrovore has the \textbf{Walker} and \textbf{Slave (Nest)} abilities. Its Bio-Flame Lance Cannon has the \textbf{Artillery} and \textbf{Reduces Cover ($-3$)} abilities.

\begin{tblr}{
    width=\textwidth,
    colspec={
        Q[c,m,1.2]
        Q[c,m,1.2]
        Q[c,m,0.7]
        Q[c,m,0.8]
        X[c,m,2.7]
        Q[c,m,1.1]
        Q[c,m,0.9]
        Q[c,m,1.2]
        Q[c,m,0.7]
    },
    row{1}={bg=black, fg=white, font=\bfseries},
    cells={valign=m},
    hlines,
    vlines
}
Movement & Save & AF & Morale & Weapon & Range & Dice & To-Hit & AP \\
15 cm & 4+ & +0 & 6 & Bio-Flame Lance Cannon & 30 cm & 2 & 3+ & 0
\end{tblr}

\unitdivider


\subsection{Mycetic Spore}

Mycetic Spores are far more than simple landing modules for the Tyranids: they are living members of the species. The Tyranids are a spacefaring race, but they hunt only upon the ground.

There are many different kinds of spores. Some affect the climate, while others alter the native fauna or flora. From a military perspective, without Mycetic Spores the Hive Fleets and Norn Queens would be forced to break through planetary defences directly and expend enormous quantities of energy merely to feed.

Mycetic Spores allow a Hive Fleet to bombard a planet with swarms of organisms and delay its final assault until the planetary defences have already been weakened.

The Tyranids' use of Mycetic Spores is comparable to a human tasting a fine wine before drinking it. The main course has not yet been served, but the meal has already begun.

Mycetic Spores have the \textbf{Deep Strike (2)}, \textbf{Mycetic Spore}, and \textbf{Transport (Special)} abilities.

\begin{tblr}{
    width=\textwidth,
    colspec={
        Q[c,m,1.2]
        Q[c,m,1.2]
        Q[c,m,0.7]
        Q[c,m,0.8]
        X[c,m,2.7]
        Q[c,m,1.1]
        Q[c,m,0.9]
        Q[c,m,1.2]
        Q[c,m,0.7]
    },
    row{1}={bg=black, fg=white, font=\bfseries},
    cells={valign=m},
    hlines,
    vlines
}
Movement & Save & AF & Morale & Weapon & Range & Dice & To-Hit & AP \\
0 cm & 3+ & +0 & -- & -- & -- & -- & -- & --
\end{tblr}

\unitdivider

\subsection{Tervigon}

The Tervigon is an enormous monster whose swollen abdomen is protected by a thick carapace. Although it possesses a formidable arsenal of bio-weapons, titanic claws capable of crushing any prey that ventures too close, and spine launchers that can strike enemies at considerable range, the greatest danger lies within its belly.

Every Tervigon is a living incubator containing numerous dormant Termagants held in a state resembling hibernation. With a mental impulse, the Tervigon can awaken its offspring at will.

An enemy attempting to engage a Tervigon may therefore find itself confronted by a swarm of Termagants emerging from its abdomen. Such a confrontation is particularly dangerous, as the Tervigon can continually reinforce the ranks of its offspring, which are driven into a battle frenzy to protect their progenitor.

A Tervigon can only transport Termagant bases.

The Tervigon has the \textbf{Slave (Devastation)}, \textbf{Regeneration (5+)}, \textbf{Attached Transport}, \textbf{Walker}, and \textbf{Transport (4 Termagants)} abilities.

\begin{tblr}{
    width=\textwidth,
    colspec={
        Q[c,m,1.2]
        Q[c,m,1.2]
        Q[c,m,0.7]
        Q[c,m,0.8]
        X[c,m,2.7]
        Q[c,m,1.1]
        Q[c,m,0.9]
        Q[c,m,1.2]
        Q[c,m,0.7]
    },
    row{1}={bg=black, fg=white, font=\bfseries},
    cells={valign=m},
    hlines,
    vlines
}
Movement & Save & AF & Morale & Weapon & Range & Dice & To-Hit & AP \\
20 cm & 2+ & +5 & Attached & Fragmentation Spines & 20 cm & 2 & 5+ & 0
\end{tblr}

\unitdivider


\subsection{Toxicrene}

The enormous monster known as the Toxicrene appears deadly through its sheer size alone. It advances across the battlefield with ponderous steps, crushing enemy warriors with its treacherous tentacles.

However, it is not the Toxicrene's immense strength that makes it one of the most feared creatures in the Hive Fleet's arsenal. Its greatest threat comes from the clouds of toxic spores expelled from the vents along its back. These spores enter the lungs of their victims and cause them to drown in their own blood.

The foul toxic emissions contain billions of embryonic microscopic organisms, each possessing a predatory awareness. When the Toxicrene releases its lethal payload, these spores invade non-Tyranid life forms and overwhelm even the most resilient immune systems, inflicting an exceptionally painful and horrific death.

The Toxicrene's emissions are so effective that neither biological protection equipment nor even Power Armour can guarantee safety from them.

The Toxicrene has the \textbf{DR (5+)}, \textbf{Walker}, and \textbf{Slave (Devastation)} abilities.

\begin{tblr}{
    width=\textwidth,
    colspec={
        Q[c,m,1.2]
        Q[c,m,1.2]
        Q[c,m,0.7]
        Q[c,m,0.8]
        X[c,m,2.7]
        Q[c,m,1.1]
        Q[c,m,0.9]
        Q[c,m,1.2]
        Q[c,m,0.7]
    },
    row{1}={bg=black, fg=white, font=\bfseries},
    cells={valign=m},
    hlines,
    vlines
}
Movement & Save & AF & Morale & Weapon & Range & Dice & To-Hit & AP \\
20 cm & 2+ & +5 & 6 & Spore Nest & -- & -- & -- & --
\end{tblr}

\unitdivider


\subsection{Virago}

The Hive Virago is a flying monstrosity used by the Tyranids to establish aerial superiority. It is perfectly adapted to aerial combat and can clear the skies of enemy aircraft regardless of their sophistication.

Hive Viragos frequently conceal their attacks behind flights of Gargoyles, hiding within the moving cloud of membranous wings until the final moment, when they break away to descend upon their prey and deliver a decisive strike.

The Virago has the \textbf{Interceptor}, \textbf{Semi-Synaptic (Devastation)}, and \textbf{Floater} abilities. Its Salivary Cannon has the \textbf{Reduces Cover ($-3$)} ability.

\begin{tblr}{
    width=\textwidth,
    colspec={
        Q[c,m,1.2]
        Q[c,m,1.2]
        Q[c,m,0.7]
        Q[c,m,0.8]
        X[c,m,2.7]
        Q[c,m,1.1]
        Q[c,m,0.9]
        Q[c,m,1.2]
        Q[c,m,0.7]
    },
    row{1}={bg=black, fg=white, font=\bfseries},
    cells={valign=m},
    hlines,
    vlines
}
Movement & Save & AF & Morale & Weapon & Range & Dice & To-Hit & AP \\
25 cm & 3+ & +6 & 5 & Salivary Cannon & 20 cm & 2 & 4+ & 0
\end{tblr}

